%=============================================================
% Annexes
%=============================================================
\chapter{Hyperparamètres des modèles ML}
\label{ann:hyperparams}

\begin{table}[H]
\centering
\caption{Hyperparamètres retenus pour chaque modèle}
\label{tab:hyperparams}
\begin{tabular}{lll}
\toprule
\textbf{Modèle} & \textbf{Paramètre} & \textbf{Valeur} \\
\midrule
\multirow{3}{*}{HistGradientBoosting}
  & \texttt{max\_iter}     & 300  \\
  & \texttt{learning\_rate} & 0,05 \\
  & \texttt{max\_depth}    & 6    \\
\midrule
\multirow{3}{*}{XGBoost}
  & \texttt{n\_estimators} & 200  \\
  & \texttt{learning\_rate} & 0,05 \\
  & \texttt{max\_depth}    & 5    \\
\midrule
\multirow{2}{*}{LogisticRegression}
  & \texttt{C}             & 1,0  \\
  & \texttt{max\_iter}     & 1000 \\
\midrule
\multirow{2}{*}{RandomForest}
  & \texttt{n\_estimators} & 200  \\
  & \texttt{max\_depth}    & 10   \\
\midrule
\multirow{2}{*}{LinearSVM}
  & \texttt{C}             & 1,0  \\
  & \texttt{max\_iter}     & 2000 \\
\bottomrule
\end{tabular}
\end{table}

\chapter{Prédictions détaillées par département}
\label{ann:dept_detail}

\begin{table}[H]
\centering
\caption{Top 3 candidats prédits par département — HistGradientBoosting}
\label{tab:dept_detail}
\begin{tabular}{lllll}
\toprule
\textbf{Département} & \textbf{1\textsuperscript{er}} & \textbf{Score} & \textbf{2\textsuperscript{e}} & \textbf{3\textsuperscript{e}} \\
\midrule
Charente            & Macron    & 41,5\% & Mélenchon & Jadot \\
Charente-Maritime   & Macron    & 64,6\% & Mélenchon & Pécresse \\
Corrèze             & Macron    & 59,2\% & Mélenchon & Jadot \\
Creuse              & Macron    & 71,8\% & Pécresse  & Lassalle \\
Deux-Sèvres         & Macron    & 75,2\% & Pécresse  & Lassalle \\
Dordogne            & Macron    & 72,5\% & Mélenchon & Pécresse \\
Gironde             & Macron    & 67,4\% & Pécresse  & Lassalle \\
Haute-Vienne        & Macron    & 52,9\% & Pécresse  & Lassalle \\
Landes              & Macron    & 73,2\% & Pécresse  & Mélenchon \\
Lot-et-Garonne      & Macron    & 67,7\% & Pécresse  & Lassalle \\
Pyrénées-Atlantiques & Macron   & 53,0\% & Mélenchon & Pécresse \\
Vienne              & Macron    & 67,0\% & Mélenchon & Pécresse \\
\bottomrule
\end{tabular}
\end{table}

\chapter{Structure des fichiers du projet}
\label{ann:structure}

\begin{verbatim}
MsprBigData/
  MSPR_Final/
    .venv/                    (Environnement Python)
    MSPR/
      01_Donnees/
        delta_lake_final/     (Parquet versionne)
        final/                (CSV finaux)
      02_Data_Engineering/
        run_prep.py           (Pipeline ETL)
      03_Data_Science/
        notebooks/
          Nouvelle_Aquitaine_Machine_Learning.ipynb
        generate_visualisations.py
        generate_tendances.py
  public/data/
    charts/                   (Images Python)
    predictions.json          (Resultats JSON)
  src/routes/
    index.tsx                 (Dashboard principal)
    visualisation.tsx         (Page visualisation)
\end{verbatim}

\chapter{Références bibliographiques}
\label{ann:refs}

Ce mémoire mobilise les ressources et travaux suivants :

\begin{itemize}
  \item INSEE — Fichiers détail du recensement de la population 2016 et 2022
    (www.insee.fr, données libres de droits sous licence Etalab)
  \item Ministère de l'Intérieur — Résultats du 1\textsuperscript{er} tour de
    l'élection présidentielle 2022 (data.gouv.fr)
  \item Pedregosa \textit{et al.}, \textit{Scikit-learn: Machine Learning in Python},
    JMLR 12, 2011
  \item Chen \& Guestrin, \textit{XGBoost: A Scalable Tree Boosting System},
    KDD 2016
  \item Ke \textit{et al.}, \textit{LightGBM: A Highly Efficient Gradient Boosting
    Decision Tree}, NeurIPS 2017
\end{itemize}
